% Chapter 3, Figure 23: flow about a circular cylinder held transverse to the airstream (scan:
% figures/ch3/fig23.png; after Schlichting). Geometry: fig23.py, in units of the cylinder radius a (1.25 cm).
% The streamlines are the plane potential flow past a doublet and a source (the cylinder with its boundary layer
% and dead water as one displacing body), fitted to the drawn streamlines; the hatched region (boundary layer and
% dead water) is bounded by the dividing streamline, curling into two lobes behind the cylinder; the velocity
% profile standing on the cylinder at B is the Blasius profile of Table 1 across the layer's thickness there.
% Points A, B, C (phi = 0, 90, 180 deg of Fig 24) and the separation point S (where the 1973 art puts it) are
% marked on the surface and tagged inside the cylinder; the curved arrows show the reversed flow in the eddies.
% The profile at B in the profile family's styles (Ch3 Figs 16, 17, 18, 25): the profile u(y) in s1 (1pt), carried
% on at U_inf (the box's right side) above the layer, as Figs 16, 18 and 25 carry theirs; the velocity lines
% headless, as printed, in s1 (profile stroke), the top one closing the box; the base line (the wall normal at B)
% and the cylinder in ink.
\documentclass[11pt]{standalone}
\usepackage{tamrfig}
\pgfplotstableread[col sep=comma]{figures/v2/ch3/fig23-marks.csv}\marks
\foreach \key in {xfront,delta,boxw,boxtop,sx,sy,xcusp,ulabx,ulaby,x0,x1} {%
  \pgfplotstablegetelem{0}{\key}\of\marks \expandafter\xdef\csname m@\key\endcsname{\pgfplotsretval}}
\makeatletter
\newcommand{\mk}[1]{\csname m@#1\endcsname}
\makeatother
% the profile family's styles (Ch3 Figs 16, 17, 18, 25)
\tikzset{
  profile/.style={draw=s1, line width=1pt},
  profile arrow/.style={draw=s1, line width=0.5pt, -{Stealth[length=3.4pt, width=2.4pt]}},
  profile stroke/.style={draw=s1, line width=0.5pt},
}
\tikzset{streamline/.style={draw=ink, line width=0.5pt},
  % a small arrowhead on a streamline, 0.95 a from its upstream end (x = -4.0 a)
  flow arrow/.style={postaction={decorate}, decoration={markings,
      mark=at position 1.19cm with {\arrow{Stealth[length=4.2pt, width=3pt]}}}},
  % an arrow over the hatching: a white casing under a thin vector
  eddy/.style={thin vec, preaction={draw=white, line width=2pt, -}}}
\begin{document}
\begin{tikzpicture}
  \begin{axis}[hide axis, x=1.25cm, y=1.25cm, disabledatascaling, anchor=origin, clip=false,
      xmin=-4.95, xmax=4.75, ymin=-2.1, ymax=2.1, every axis plot/.append style={mark=none}]
    % streamlines: four above and four below, the stagnation streamline ahead and the wake's centre line
    \pgfplotsinvokeforeach{1,2,3,4}{
      \addplot[streamline, flow arrow] table[x=x, y=s#1, col sep=comma] {figures/v2/ch3/fig23-stream.csv};
      \addplot[streamline, flow arrow] table[x=x, y expr=-\thisrow{s#1}, col sep=comma] {figures/v2/ch3/fig23-stream.csv};
    }
    \draw[streamline, flow arrow] (\mk{x0},0) -- (\mk{xfront},0);
    \draw[streamline] (\mk{xcusp},0) -- (\mk{x1},0);
    \node[fill=white, inner sep=1pt] at (\mk{ulabx},\mk{ulaby}) {$U_\infty$};
    % boundary layer and dead water, then the cylinder over it
    \addplot[outline, hatch] table[x=x, y=y, col sep=comma] {figures/v2/ch3/fig23-wake.csv} -- cycle;
    \draw[outline, fill=white] (0,0) circle[radius=1];
    \draw[centerline] (-0.56,0) -- (0.56,0);
    \draw[centerline] (0,-0.86) -- (0,0.56);
    % the separated layer's outer flow, downstream from beside S (as printed), and the reversed flow round each
    % eddy (clockwise above, mirrored below)
    \pgfplotsinvokeforeach{1,-1}{
      \draw[eddy] (0.86,{#1*0.84}) -- (1.16,{#1*0.71});
      \draw[eddy] (1.19,{#1*0.38}) .. controls (1.22,{#1*0.58}) and (1.44,{#1*0.62}) .. (1.66,{#1*0.46});
    }
    % the velocity profile at B: a white box standing on the cylinder (clearing the hatching and the streamlines),
    % its velocity lines, the top one closing the box, the profile (Table 1, then U_inf to the top), the base line
    \addplot[fill=white, draw=none] table[x=x, y=y, col sep=comma] {figures/v2/ch3/fig23-box.csv} -- cycle;
    \addplot[profile stroke, quiver={u=\thisrow{len}, v=0}, -]
      table[x=x, y=y, col sep=comma] {figures/v2/ch3/fig23-ticks.csv};
    \draw[profile stroke] (0,\mk{boxtop}) -- (\mk{boxw},\mk{boxtop});
    \addplot[profile, line join=round] table[x=x, y=y, col sep=comma] {figures/v2/ch3/fig23-profile.csv}
      -- (\mk{boxw},\mk{boxtop});
    \draw[outline] (0,1) -- (0,\mk{boxtop});
    % the points A, B, C and S
    \node[point] at (-1,0) {};  \node[curve tag] at (-0.79,0) {A};
    \node[point] at (0,1) {};   \node[curve tag] at (0,0.79) {B};
    \node[point] at (1,0) {};   \node[curve tag] at (0.79,0) {C};
    \node[point] at (\mk{sx},\mk{sy}) {};  \node[curve tag] at ({0.79*\mk{sx}},{0.79*\mk{sy}}) {S};
  \end{axis}
\end{tikzpicture}
\end{document}
